\documentclass[xcolor=table,dvipsnames,table]{beamer}
\mode<presentation>
\usetheme{boxes}
\setbeamertemplate{navigation symbols}{}
% http://www.latex-community.org/forum/viewtopic.php?f=4&t=6694
\setbeamertemplate{navigation symbols}{\raisebox{5pt}{\makebox[\paperwidth]{\hfill\makebox[10pt]{\scriptsize\insertframenumber\vspace{1ex}}}}}
%\setbeamertemplate{footline}[frame number]
\setbeamertemplate{blocks}[shadow=false]
%\setbeamercolor*{block title}{fg=structure,bg=RoyalBlue!10}
\setbeamercolor*{block title example}{fg=structure,bg=RoyalBlue!10}
%\setbeamercolor*{block title example}{fg=BrickRed,bg=Goldenrod!10}
\setbeamercolor*{block title alerted}{fg=white,bg=black}
\addtobeamertemplate{block begin}{\pgfsetfillopacity{0.8}}{\pgfsetfillopacity{1}}
%\rowcolors{0}{RoyalBlue!20}{RoyalBlue!5}
\setbeamertemplate{caption}{\raggedright\insertcaption\par}

%\DeclareGraphicsRule{*}{mps}{*}{}

\usepackage{latexsym}
\usepackage{hyperref}
\usepackage{tikz}
\usetikzlibrary{calc,shapes,arrows,shadows,shapes.callouts,shapes.arrows,chains,positioning,trees}
\usepackage{solution}
\usepackage{calc}
\usepackage{pifont}
\usepackage{algorithmic}
\usepackage{pdfcomment}
\usepackage{color}

\newcommand{\cmark}{\ding{51}}
\newcommand{\xmark}{\ding{55}}

\newcommand{\highlight}[1]{{\color{blue}{#1}}}
\newcommand{\mycite}[1]{{\color{darkgray}{\footnotesize [#1]}}}

\DeclareMathOperator*{\argmin}{arg\,min}
\DeclareMathOperator*{\argmax}{arg\,max}
\DeclareMathOperator{\sign}{sign}
\DeclareMathOperator{\cnt}{Count}

\newcounter{mycallout}

\newcommand{\callouts}[3]{%
  \stepcounter{mycallout}
  \tikz[remember picture,baseline]{\node[anchor=base,inner sep=0,outer sep=0]%
    (\themycallout) {\colorbox{#1!20}{#3}};\pause\node[overlay,rectangle callout,%
    callout relative pointer={(0cm,0.5cm)},fill=#1!20] at ($(\themycallout.south)+(-0cm,-0.7cm)$){#2};}%
    }%

\raggedright

\newcount\lecturecount
\lecturecount=0
\AtBeginLecture{%
    \advance\lecturecount by 1
    \date{}
    \begin{frame}
    \begin{center}
    \titlepage
    \ifnum\lecturecount=1
    Part \the\lecturecount: \insertlecture
    \else
    Part \the\lecturecount: \insertlecture
    \fi
    \end{center}
    \end{frame}
}

\addtobeamertemplate{block begin}{\setlength\abovedisplayskip{0pt}}

%\newcommand{\example}[1]{{\color{BrickRed!50}{#1}}}
\newcommand{\maths}[1]{{\color{RoyalBlue!50}{#1}}}
\newcommand{\reference}[1]{{\color{RoyalBlue!30}\tiny [from #1]}}
\newcommand{\koehnref}{\reference{\href{http://www.statmt.org/book}{P.Koehn SMT book slides}}}


\begin{document}

\title{\color{MidnightBlue}Natural Language Processing}

\author{Anoop Sarkar \\ {\color{RoyalBlue!70}{\href{http://anoopsarkar.github.io/nlp-class}{anoopsarkar.github.io/nlp-class}}}}
\institute{\color{BrickRed}Simon Fraser University}
%\date{}
     
{
\addtocounter{framenumber}{-1}
\begin{frame}
\begin{center}
\vspace{8mm}
\includegraphics[scale=0.35]{figures/natlang-cky-logo}
\end{center}
\titlepage
\end{frame}
}



\lecture{Classification tasks in NLP}{}

\section{Classification tasks in NLP}
\frame{\tableofcontents[currentsection]}

\begin{frame}
\frametitle{Prepositional Phrases}
\begin{itemize}[<+->]
\item noun attach: {\em I bought the shirt with pockets}
\item verb attach: {\em I washed the shirt with soap}
\item As in the case of other attachment decisions in parsing: it depends on the meaning of the entire sentence -- needs world knowledge, etc.
\item Maybe there is a simpler solution: we can attempt to solve it using heuristics or associations between words
\end{itemize}

\end{frame}

\begin{frame}
\frametitle{Ambiguity Resolution: Prepositional Phrases in English}
  \begin{itemize}[<+->]
  \item Learning Prepositional Phrase Attachment: Annotated Data
\begin{tabular}{|cccc|c|}
\hline
$v$    &     $n_1$   &      $p$ & $n_2$ &       Attachment\\
\hline
join   &   board    &   as &  director & V \\
is       & chairman  &  of  & N.V.    & N \\
using  &   crocidolite & in  & filters & V \\
bring  &   attention  & to  & problem &  V \\ 
is      &  asbestos   & in &  products & N \\
making &   paper    &   for & filters & N \\
including & three     &  with & cancer &  N \\
$\vdots$ & $\vdots$ & $\vdots$ & $\vdots$ & $\vdots$ \\
\hline
\end{tabular}
  \end{itemize}

\end{frame}

\begin{frame}
\frametitle{Prepositional Phrase Attachment}
\begin{tabular}{|l|c|}  \hline
Method & Accuracy \\ \hline
Always noun attachment & 59.0 \\
Most likely for each preposition & 72.2 \\
Average Human (4 head words only) & 88.2 \\
Average Human (whole sentence) & 93.2 \\  \hline
\end{tabular}

\end{frame}

\begin{frame}
\frametitle{Back-off Smoothing}
\begin{itemize}[<+->]
\item Random variable $a$ represents attachment. 
\item $a = n_1$ or $a = v$ (two-class classification)
\item We want to compute probability of noun attachment: $p(a = n_1 \mid v, n_1, p, n_2)$. 
\item Probability of verb attachment is $1 - p(a = n_1 \mid v, n_1, p, n_2)$.
\end{itemize}
\end{frame}

\begin{frame}
\frametitle{Back-off Smoothing}
\begin{enumerate}
\item<1-> If $f(v,n_1,p,n_2) > 0$ and $\hat{p} \neq 0.5$
{\small
\[ \hat{p}(a_{n_1} \mid  v,n_1,p,n_2) = \frac{ f(a_{n_1},v,n_1,p,n_2) }{ f(v,n_1,p,n_2)
} \]
}
\item<2-> Else if $f(v,n_1,p) + f(v,p,n_2) + f(n_1,p,n_2) > 0$ \\
and $\hat{p} \neq 0.5$
{\small
\[ \hat{p}(a_{n_1} \mid  v,n_1,p,n_2) = \frac{ f(a_{n_1},v,n_1,p) + f(a_{n_1},v,p,n_2) +
  f(a_{n_1},n_1,p,n_2) }{ f(v,n_1,p) + f(v,p,n_2) + f(n_1,p,n_2) } \]
  }
\item<3-> Else if $f(v,p) + f(n_1,p) + f(p,n_2) > 0$
{\small
\[ \hat{p}(a_{n_1} \mid  v,n_1,p,n_2) = \frac{ f(a_{n_1},v,p) + f(a_{n_1},n_1,p) +
  f(a_{n_1},p,n_2) }{ f(v,p) + f(n_1,p) + f(p,n_2) } \]
}
\item<4-> Else if $f(p) > 0$ (try choosing attachment based on preposition alone)
{\small
\[ \hat{p}(a_{n_1} \mid  v,n_1,p,n_2) = \frac{ f(a_{n_1},p) }{ f(p) } \]
}\item<5-> Else {\small
\( \hat{p}(a_{n_1} \mid  v,n_1,p,n_2) = 1.0 \)
}
\end{enumerate}
\end{frame}

\begin{frame}
\frametitle{Prepositional Phrase Attachment: Results}
  \begin{itemize}[<+->]
  \item {\bf Results (Collins and Brooks 1995)}: 84.5\% accuracy \\
  with the use of some limited word classes for dates, numbers, etc.
  \item {\bf Toutanova, Manning, and Ng, 2004}: \\
use sophisticated smoothing model for PP attachment\\
86.18\% with words \& stems; with word classes: 87.54\%
  \item {\bf Merlo, Crocker and Berthouzoz, 1997}:\\
 test on multiple PPs, generalize disambiguation of 1 PP to 2-3 PPs\\
%14 structures possible for 3PPs assuming a single verb: all 14 are attested in the Treebank\\
%same model as CB95; but generalized to dealing with upto 3PPs\\
1PP: 84.3\% \ \ 2PP: 69.6\% \ \ 3PP: 43.6\% \\
{\bf Note that this is still not the real problem faced in parsing natural language}
  \end{itemize}
\end{frame}

\lecture{Probabilistic Classifiers}{}

\section{Naive Bayes Classifier}
\frame{\tableofcontents[currentsection]}

\begin{frame}
\frametitle{Naive Bayes Classifier}
\begin{itemize}[<+->]
\item \textbf{x} is the input that can be represented as $d$ independent features $f_j$, $1 \leq j \leq d$
\item $y$ is the output classification
\item $P(y \mid {\bf x}) = \frac{P(y) \cdot P({\bf x} \mid y)}{P({\bf x})}$ (Bayes Rule)
\item $P({\bf x} \mid y) = \prod_{j=1}^{d} P(f_j  \mid y)$
\item $P(y \mid {\bf x}) = P(y) \cdot  \prod_{j=1}^{d} P(f_j  \mid y)$
\end{itemize}
\end{frame}

\section{Log linear models}
\frame{\tableofcontents[currentsection]}

\begin{frame}
\frametitle{Log linear model}
\begin{itemize}[<+->]
\item Let there be $m$ features, $f_k(\textbf{x}, y)$ for $k = 1, \ldots, m$
\item Define a parameter vector $\textbf{w} \in \mathbb{R}^m$
\item Each $(\textbf{x}, y)$ pair is mapped to score:
\[ s(\textbf{x}, y) = \sum_k w_k \cdot f_k(\textbf{x}, y) \]
\item Using inner product notation:
\begin{eqnarray*}
\textbf{w} \cdot \textbf{f}(\textbf{x}, y) & = & \sum_k w_k \cdot f_k(\textbf{x}, y) \\
s(\textbf{x}, y) & = & \textbf{w} \cdot \textbf{f}(\textbf{x}, y) 
\end{eqnarray*}
\item To get a probability from the score: Renormalize! 
\[ \Pr(y \mid \textbf{x}, \textbf{w}) = \frac{exp\left(s(\textbf{x}, y)\right)}{\sum_{y'} exp\left(s(\textbf{x}, y')\right)} \]
\end{itemize}
\end{frame}

\begin{frame}
\frametitle{Log linear model}
\begin{itemize}[<+->]
\item The name `log-linear model' comes from:
\[ \log \Pr(y \mid \textbf{x}, \textbf{w}) = \underbrace{\textbf{w} \cdot \textbf{f}(\textbf{x}, y)}_{\textrm{linear term}} - \underbrace{\log \sum_{y'} exp\left( \textbf{w} \cdot \textbf{f}(\textbf{x}, y') \right)}_{\textrm{normalization term}} \]
\item Once the weights are learned, we can perform predictions using these features.
\item The goal: to find $\textbf{w}$ that maximizes the log likelihood $L(\textbf{w})$ of the labeled training set containing $(\textbf{x}_i, y_i)$ for $i = 1 \ldots n$
\begin{eqnarray*}
L(\textbf{w}) &=& \sum_{i} \log \Pr(y_i \mid \textbf{x}_i, \textbf{w}) \\
&=& \sum_i \textbf{w} \cdot \textbf{f}(\textbf{x}_i, y_i) - \sum_i \log \sum_{y'} exp\left( \textbf{w} \cdot \textbf{f}(\textbf{x}_i, y') \right) 
\end{eqnarray*}
\end{itemize}
\end{frame}

\begin{frame}
\frametitle{Log linear model}
\begin{itemize}[<+->]
\item Maximize:
\begin{eqnarray*}
L(\textbf{w}) &=& \sum_i \textbf{w} \cdot \textbf{f}(\textbf{x}_i, y_i) - \sum_i \log \sum_{y'} exp\left( \textbf{w} \cdot \textbf{f}(\textbf{x}_i, y') \right) 
\end{eqnarray*}
\item Calculate gradient:
\begin{eqnarray*}
\lefteqn{\left. \frac{d L(\textbf{w})}{d \textbf{w}} \right|_{\textbf{w}} } \\
&=& \sum_i \textbf{f}(\textbf{x}_i, y_i) - \sum_i \frac{1}{\sum_{y''} exp\left(\textbf{w} \cdot \textbf{f}(x_i, y'')\right)} \\
&& \sum_{y'} \textbf{f}(\textbf{x}_i, y')  \cdot exp\left( \textbf{w} \cdot \textbf{f}(\textbf{x}_i, y') \right) \\
&=& \sum_i \textbf{f}(\textbf{x}_i, y_i) - \sum_i \sum_{y'} \textbf{f}(\textbf{x}_i, y') \frac{exp\left( \textbf{w} \cdot \textbf{f}(\textbf{x}_i, y') \right)}{\sum_{y''} exp\left(\textbf{w} \cdot \textbf{f}(x_i, y'')\right)} \\
&=& \underbrace{\sum_i \textbf{f}(\textbf{x}_i, y_i)}_{\textrm{Observed counts}} - \underbrace{\sum_i \sum_{y'} \textbf{f}(\textbf{x}_i, y') \Pr(y' \mid \textbf{x}_i, \textbf{w})}_{\textrm{Expected counts}}
\end{eqnarray*}
\end{itemize}
\end{frame}

\begin{frame}
\frametitle{Log linear model}
\begin{itemize}[<+->]
\item Init: $\textbf{w}^{(0)} = \textbf{0}$
\item $t \leftarrow 0$
\item Iterate until convergence:
\begin{itemize}[<+->]
\item Calculate: $\Delta = \left. \frac{d L(\textbf{w})}{d \textbf{w}}  \right|_{\textbf{w} = \textbf{w}^{(t)}}$
\item Find $\beta^\ast = \arg\max_\beta L(\textbf{w}^{(t)} + \beta \Delta)$
\item Set $\textbf{w}^{(t+1)} \leftarrow \textbf{w}^{(t)} + \beta^\ast \Delta$
\end{itemize}
\end{itemize}
\end{frame}

\begin{frame}
\frametitle{Learning the weights: $\textbf{w}$: Generalized Iterative
  Scaling}
\begin{tabbing}
$f^\# = max_{x,y} \sum_{j} f_j(x, y)$ \\
(the maximum possible feature value; needed for scaling) \pause \\
Initialize $\textbf{w}^{(0)}$ \pause \\
For \= each iteration $t$ \\
\> \texttt{expected[j]} $\leftarrow$ 0 for j = 1 .. \# of features \pause \\
\> For \= i = 1 to $\mid \textrm{training data} \mid$ \\
\>     \> For \= each feature $f_j$ \\
\>     \>     \> \texttt{expected[j]} += $f_j(x_i, y_i) \cdot P(y_i \mid x_i, \textbf{w}^{(t)})$ \pause \\
\> For \= each feature $f_j(x,y)$ \\
\>     \> \texttt{observed[j]} = $f_j(x, y) \cdot \frac{c(x,y)}{\mid \textrm{training data} \mid}$ \pause \\ 
\> For \= each feature $f_j(x,y)$ \\
\>     \> $w_j^{(t+1)} \leftarrow w_j^{(t)} \cdot \sqrt[f^\#]{\frac{\texttt{observed[j]}}{\texttt{expected[j]}}}$ 
\end{tabbing}
\par\noindent
\small{cf. Goodman, NIPS '01}
\end{frame} 

\lecture{Linear models for Tagging}{}

\section{Tagging tasks in NLP}
\frame{\tableofcontents[currentsection]}
\newcommand{\postag}[1]{{\color{red}/#1}}
\newcommand{\nertag}[1]{{\color{blue}/#1}}

\begin{frame}
\frametitle{Tagging Tasks}
\begin{block}{Tagged Sequences}
a b e e a f h j $\Rightarrow$ a\postag{Y} b\postag{Z} e\postag{Y} e\postag{Y} a\postag{Z} f\postag{X} h\postag{Z} j\postag{Y}
\end{block}
\pause
\begin{alertblock}{Example 1: Part-of-speech tagging}
Profits\postag{N} soared\postag{V} at\postag{P} Boeing\postag{N} Co.\postag{N} ,\postag{,} easily\postag{ADV} topping\postag{V} forecasts\postag{N} on\postag{P} Wall\postag{N} Street\postag{N} ,\postag{,} as\postag{P} their\postag{POSS} CEO\postag{N} Alan\postag{N} Mulally\postag{N} announced\postag{V} first\postag{ADJ} quarter\postag{N} results\postag{N} .\postag{.}
\end{alertblock}
\pause
\begin{alertblock}{Example 2: Named Entity Recognition}
Profits\postag{O} soared\postag{O} at\postag{O} Boeing\nertag{B-CO} Co.\nertag{I-CO} ,\postag{O} easily\postag{O} topping\postag{O} forecasts\postag{O} on\postag{O} Wall\nertag{B-LOC} Street\nertag{I-LOC} ,\postag{O} as\postag{O} their\postag{O} CEO\postag{O} Alan\nertag{B-PER} Mulally\nertag{I-PER} announced\postag{O} first\postag{O} quarter\postag{O} results\postag{O} .\postag{O}
\end{alertblock}
\end{frame}


\begin{frame}
\frametitle{Notation for Tagging Tasks}
\begin{itemize}[<+->]
\item Set of possible input words: ${\cal V}$
\item Set of possible tags: ${\cal T}$
\item Word sequence: $x_{[1:n]} = [x_1, \ldots, x_n]$
\item Tag sequence: $t_{[1:n]} = [t_1, \ldots, t_n]$
\item Training data is $N$ tagged sentences, the $i^{th}$ sentence has length $n_i$:
\[ (x_{[1:n]}^{(i)}, t_{[1:n]}^{(i)}) \textrm{ for } i = 1, \ldots, n \]
\end{itemize}
\end{frame}

\begin{frame}
\frametitle{Independence Assumptions for Tagging}
\begin{block}{Chain Rule}
\[ P( t_{[1:n]} \mid x_{[1:n]}) = \pause \prod_{j=1}^n P(t_j \mid t_{j-1}, \ldots, t_1, x_{[1:n]}, j) \]
\end{block}
\pause
\begin{block}{Make independence assumptions}
\[ P( t_{[1:n]} \mid x_{[1:n]}) \approx \pause \prod_{j=1}^n P(t_j \mid t_{j-1}, x_{[1:n]}, j) \]
\centering {\small $j$ is the word being tagged. }\\
\centering {\color{red}\small We model the conditional probability directly: no Bayes Rule here. }
\end{block}
\pause
\begin{block}{Questions}
\begin{itemize}
\item Split up $P(t_j \mid t_{j-1}, x_{[1:n]}, j)$ into parameters?
\item How to find $\arg\max_{t_{[1:n]}} P( t_{[1:n]} \mid x_{[1:n]})$?
\end{itemize}
\end{block}
\end{frame}

\section{Log-linear models for Tagging}
\frame{\tableofcontents[currentsection]}

\begin{frame}
\frametitle{Representation: finding the right parameters}
\begin{block}{Problem: Predict ?? using context, $P(?? \mid \texttt{context})$ }
Profits\postag{N} soared\postag{V} at\postag{P} Boeing\postag{??} Co. , easily topping forecasts on Wall Street , as their CEO Alan Mulally announced first quarter results .
\end{block}
\pause
\begin{block}{Representation: history}
\begin{itemize}
\item A history is a 3-tuple: $(t_{-1}, x_{[1:n]}, i)$
\item $t_{-1}$ is the previous tag (we are assuming a bigram model)
\item $x_{[1:n]}$ are the $n$ words in the input
\item $i$ is the index of the word being tagged
\item For example, for $x_4$ = Boeing:
    \begin{itemize}
    \item $t_{-1}$ = P
    \item $x_{[1:n]}$ = (Profits, soared, ..., results, .)
    \item $i$ = 4
    \end{itemize}
\end{itemize}
\end{block}
\end{frame}

\newcommand{\featurefunction}[3]{
\[ 
f_{#1}(h,t) = \left\{
\begin{array}{ll}
1 & \textrm{if #2 and $t$ = #3 } \\
0 & \textrm{otherwise}
\end{array}
\right.
\]
}

\begin{frame}
\frametitle{Feature-vectors over history-tag pairs}
\begin{block}{Take a history, tag pair $(h,t)$}
$f_k(h, t)$ for $k = 1, \ldots, m$ are \textbf{feature functions} representing the tagging decision.
\end{block}
\pause
\begin{block}{Example: Part-of-speech tagging [Ratnaparkhi 1996]}
\centering
\featurefunction{100}{current word $x_i$ is \texttt{Boeing}}{\texttt{N}}
\featurefunction{101}{$t_{-1}$ is \texttt{P}}{\texttt{N}}
\end{block}
\end{frame}

\begin{frame}
\frametitle{Log linear model for Tagging}
\begin{itemize}[<+->]
\item Let there be $m$ features, $f_k(\textbf{x}, \textbf{y})$ for $k = 1, \ldots, m$
\item $\textbf{x} = x_{[1:n]}$ and $\textbf{y} = t_{[1:n]}$
\item Define a parameter vector $\textbf{w} \in \mathbb{R}^m$
\item Each $(\textbf{x}, \textbf{y})$ pair is mapped to score:
\[ s(\textbf{x}, \textbf{y}) = \sum_k w_k \cdot f_k(\textbf{x}, \textbf{y}) \]
\item Using inner product notation:
\begin{eqnarray*}
\textbf{w} \cdot \textbf{f}(\textbf{x}, \textbf{y}) & = & \sum_k w_k \cdot f_k(\textbf{x}, \textbf{y}) \\
s(\textbf{x}, \textbf{y}) & = & \textbf{w} \cdot \textbf{f}(\textbf{x}, \textbf{y}) 
\end{eqnarray*}
\item To get a probability from the score: Renormalize! 
\[ \Pr(\textbf{y} \mid \textbf{x}, \textbf{w}) = \frac{exp\left(s(\textbf{x}, \textbf{y})\right)}{\sum_{\textbf{y'}} exp\left(s(\textbf{x}, \textbf{y'})\right)} \]
\end{itemize}
\end{frame}

\begin{frame}
\frametitle{Feature functions for Tagging}
\begin{block}{Problem}
\begin{itemize}
\item We have defined a log-linear model using feature functions: $\textbf{f(x,y)}$ 
\item We have defined parameters using a history $h$ so feature functions are: $\textbf{f}(h, t)$
\end{itemize}
\end{block}
\end{frame}

\begin{frame}
\frametitle{Locally normalized log-linear taggers}
\begin{block}{Conditional Distribution over history, tag pair $(h,t)$}
\[ \log \Pr(t \mid h) = \textbf{w} \cdot \textbf{f}(h, t) - log \sum_{t'} \exp \left( \textbf{w} \cdot \textbf{f}(h, t') \right) \]
\begin{itemize}
\item $\textbf{f}(h, t)$ is a vector of feature functions
\item $\textbf{w}$ is the weight vector
\end{itemize}
\end{block}
\pause
\begin{block}{Local normalization for tagging}
\begin{itemize}
\item Word sequence: $x_{[1:n]}$ and tag sequence: $t_{[1:n]}$
\item Histories $h_i = (t_{i-1}, x_{[1:n]}, i)$ 
\end{itemize}
\[ \log \Pr( t_{[1:n]} \mid x_{[1:n]} ) = \sum_{i=1}^n \log \Pr(t_i \mid h_i) \]
\end{block}
\end{frame}


\begin{frame}
\frametitle{Globally normalized log-linear taggers}
\begin{block}{Global feature function $\mathbf{\Phi}(\textbf{x}, \textbf{y})$}
\begin{itemize}
\item Word sequence: $\textbf{x} = x_{[1:n]}$ and tag sequence: $\textbf{y} = t_{[1:n]}$
\item From \textit{local} histories $h_i = (t_{i-1}, x_{[1:n]}, i)$ to global $\Phi$ values:
\end{itemize}
\[ \Phi_k(x_{[1:n]}, t_{[1:n]}) = \sum_{i=1}^n f_k(h_i, t_i) \]
\begin{itemize}
\item $\mathbf{\Phi}(\textbf{x}, \textbf{y}) = (\Phi_1, \Phi_2, \ldots, \Phi_m)$ is a \textit{global} feature vector
\item $\textbf{w}$ is the weight vector for $\mathbf{\Phi}$
\end{itemize}
\end{block}
\pause
\begin{block}{Global normalization for tagging}
\[ \log \Pr(\textbf{y} \mid \textbf{x}, \textbf{w}) = \textbf{w} \cdot \mathbf{\Phi}(\textbf{x}, \textbf{y}) - \log \sum_{\textbf{y'}} exp\left( \textbf{w} \cdot \mathbf{\Phi}(\textbf{x}, \textbf{y'}) \right) \]
\end{block}
\end{frame}

\begin{frame}
\frametitle{Conditional Random Field}
\begin{block}{Global normalization for tagging}
\[ \log \Pr(\textbf{y} \mid \textbf{x}, \textbf{w}) = \textbf{w} \cdot \mathbf{\Phi}(\textbf{x}, \textbf{y}) - \log \sum_{\textbf{y'}} exp\left( \textbf{w} \cdot \mathbf{\Phi}(\textbf{x}, \textbf{y'}) \right) \]
\begin{itemize}
\item This model is also called a conditional random field (CRF)
\end{itemize}
\end{block}
\pause
\begin{block}{Algorithms for training and decoding}
\begin{itemize}
\item Global normalization could be expensive: requires sum over exponentially many terms $\textbf{y'}$
\item Finding $\arg\max_{\textbf{y}} \log \Pr(\textbf{y} \mid \textbf{x})$ can be accomplished using the Viterbi algorithm.
\item Training: finding the weight vector $\textbf{w}$ can be done using a variant of the Forward algorithm.
\end{itemize}
\end{block}
\end{frame}


\section*{Acknowledgements}

\begin{frame}
\centering
\begin{alertblock}{Acknowledgements}
Many slides borrowed or inspired from lecture notes by Michael Collins, Chris Dyer, Kevin Knight, Philipp Koehn, Adam Lopez, Graham Neubig and Luke Zettlemoyer from their NLP course materials. 

\bigskip

All mistakes are my own.

\bigskip

A big thank you to all the students who read through these notes and helped me improve them.

\end{alertblock}
\end{frame}



\end{document}

\begin{frame}
\frametitle{Maximum Entropy}
\begin{itemize}[<+->]
\item The log-linear model has an interpretation as a {\it maximum entropy} model.
\item For observed events, maximize likelihood. For unobserved events, from all
consistent models pick the one with maximum entropy.
\item The maximum entropy principle: related to Occam's razor and
  other similar justifications for scientific inquiry
\item Make the minimum possible assumptions about unseen data
\item Also: Laplace's {\em Principle of Insufficient Reason}: when one
  has no information to distinguish between the probability of two
  events, the best strategy is to consider them equally likely
\end{itemize}
\end{frame}

\begin{frame}
\frametitle{Logistic Regression}
\begin{itemize}
\item models effects of explanatory variables on binary valued
  variable
\item observations ${\bf x} = \{x_1, \ldots, x_j\}$ with success given
  by $q({\bf x})$: \[ q({\bf x}) = \frac{e^{g({\bf x})}}{1 + e^{g({\bf
  x})}} \] and \[ g({\bf x}) = \beta_0 + \sum_{j=1}^{k} \beta_j x_j \]
\end{itemize} 
\end{frame} 

\begin{frame}
\frametitle{Logistic Regression}
\begin{itemize}
\item probability that observations lead to success, or $p(a=1 \mid
  b)$:
 \[ p(a=1 \mid b) = \frac{e^{g(b)}}{1 + e^{g(b)}} \] where
 \[ g(b) = \beta_0 f_0(1,b) + \sum_{j=1}^{k} \beta_j f_j(1,b) \]
\item $\beta_j = \textrm{log} \alpha_j$, $f_0(1,b) = 1$ and $f_j(1,b) = x_j$
\end{itemize} 
\end{frame} 
